\section{Local dynamics and full-state matching}\label{sec:matching}
A related matching picture appeared earlier with the candidate constant in
Douglas's v1.0.0 archive~\cite{Douglas2026Zenodo}. The purpose of this section
is to formulate and prove the bridge between the intrinsic full-state matching
set of Definition~\ref{def:matching} and the carriers used here, together with
the uniqueness architecture developed below.

\subsection{Outer fixed points and hyperbolicity}
Use the outer coordinate
\begin{equation}\label{eq:sigma}
 t=a(s),\qquad \sigma(t)=\frac{4t^4-5t^2+2}{4t^2-1},\qquad
 T=a(I)=(a(L),a(U)].
\end{equation}
The equation \(s=\sigma(t)\) is
\(4t^4-(4s+5)t^2+s+2=0\); its larger squared root is \eqref{eq:a}.
On this branch \(\sqrt3/2<t<1\), and
\begin{equation}\label{eq:sigmaprime}
 \sigma'(t)=\frac{2t(4t^2-3)(4t^2+1)}{(4t^2-1)^2}>0.
\end{equation}
Exact comparison through \(\sigma\) gives
\(0.8782729451808<a(s)<0.87827295218\) on all of \(I\), including its
upper endpoint. This is an enclosure of the full initial domain, not a
restriction to a neighborhood of a suspected root.

\begin{premise}[Real-analytic hyperbolic local manifolds]\label{prem:stable}
A real-analytic local diffeomorphism at a hyperbolic fixed point has analytic
local stable and unstable manifolds tangent to the corresponding invariant
subspaces. In sufficiently small admissible neighborhoods their local
convergence sets capture the appropriate forward and backward convergent
orbits; the manifold germs are unique. We use the standard finite-dimensional
local theorem, with analytic regularity as in~\cite{CabreFontichLlave2003},
Theorem 1.1 and the stable/unstable specialization, and the usual local
capture formulation~\cite{Shub1987}.
\end{premise}

\begin{lemma}[The intrinsic local tails]\label{lem:tails}
For every \(s\in I\),
\[
 p_\pm(s)=(\pm a(s),\pm a(s),r(\pm a(s)))
\]
are valid fixed points. The positive fixed point has one unstable and two
stable directions; the negative fixed point has one stable and two unstable
directions. Their splittings are uniform over the closed parameter hull.
The decreasing backward convergence set near \(p_+\) and decreasing forward
convergence set near \(p_-\) are the corresponding one-dimensional analytic
local tail pieces.
\end{lemma}
\begin{proof}
Substitution in \eqref{eq:F} proves the fixed-point equations and
\(r(t)r(-t)=1\). Further,
\[
 r(t)^2-1=\frac{2t(4t^2-3)}{(1-t)(2t+1)^2}>0.
\]
At \(p_+\), one eigenvalue is \(\nu=1/r(t)^2\in(0,1)\). The other two
are the zeros of the sign-equivalent characteristic factor
\[
 (4t^3-3t+1)\xi^2-4t^2(4t^2-1)\xi-(4t^3-3t-1).
\]
Its leading and constant coefficients are positive, its root sum is
positive, and its value at one is
\(-2(2t^2-1)(4t^2+1)<0\). Hence its roots satisfy
\(0<\alpha<1<\mu\). The spectrum at \(p_-\) is reciprocal.
The determinants at the two points are respectively
\[
 \frac{(t-1)^2(2t+1)^4}{(t+1)^2(2t-1)^4},\qquad
 \frac{(t+1)^2(2t-1)^4}{(t-1)^2(2t+1)^4},
\]
both nonzero. These are exact symbolic identities, specified in
Appendix~\ref{app:symbolic}. Positive radicands and ratios make \(F_s\)
analytic near the points; its nonzero determinant gives local invertibility.
Compactness of \([a(L),a(U)]\) and continuous eigenvalue branches, none on
the unit circle, give uniform hyperbolicity. External
premise~\ref{prem:stable} now supplies the stated manifolds and local capture.
\end{proof}

\subsection{Intrinsic definitions and the crossing convention}
Choose sufficiently small admissible open neighborhoods \(N_\pm\) of the
fixed points. Define \(W^u_{s,\mathrm{loc}}\) to consist of states in
\(N_+\) admitting a valid backward chain wholly in \(N_+\) which converges
to \(p_+\); define \(W^s_{s,\mathrm{loc}}\) using a valid forward chain
wholly in \(N_-\) converging to \(p_-\). Superscript \(\downarrow\)
restricts these sets to decreasing tails. In particular it excludes the
fixed point itself. These two tail sets are defined independently.

\begin{definition}[Full-state matching]\label{def:matching}
For \(s\in I\), define
\begin{align*}
 \mathcal U_s=\{X:\;&F_s^m(\xi)=X\text{ along a valid decreasing chain},\\
 &m\ge0,\ \xi\in W^{u,\downarrow}_{s,\mathrm{loc}}\},\\
 \mathcal V_s=\{X:\;&F_s^n(X)=\eta\text{ along a valid decreasing chain},\\
 &n\ge0,\ \eta\in W^{s,\downarrow}_{s,\mathrm{loc}}\}.
\end{align*}
The integers and the two local tail points are independent. Set
\begin{equation}\label{eq:matching}
 \Sigma=\{(x,y,u)\in\mathcal D:x\ge0>y\},\qquad
 \Ms=\Sigma\cap\mathcal U_s\cap\mathcal V_s,
\end{equation}
and \(\mathcal M=\{(s,X):s\in I,\ X\in\Ms\}\).
All three state coordinates must agree at the join. Equality of the two
labels alone is insufficient. The half-open crossing set chooses a
representative of a discrete orbit; smooth transversality is not assumed.
\end{definition}

The global continuations are independent of the sufficiently small admissible
neighborhoods. Indeed a backward tail tending to \(p_+\) eventually lies
entirely in any second such neighborhood. An earlier point on that same tail
represents the same global state by finitely many valid decreasing steps.
Interchanging the neighborhoods gives equality of the source continuations;
the forward argument gives equality of the target continuations.

\begin{proposition}[Matching soundness and carrier completeness]\label{prop:bridge}
For every \(s\in I\),
\begin{equation}\label{eq:bridge}
 \mathcal C_s\ne\varnothing\quad\Longleftrightarrow\quad
 \mathcal M_s\ne\varnothing.
\end{equation}
Every carrier has a unique index \(k\) with \(c_k\ge0>c_{k+1}\), and its
state there belongs to \(\Ms\). Every match reconstructs an admissible
normalized positive carrier.
\end{proposition}
\begin{proof}
Given a match, concatenate its backward and forward tails at their identical
three-coordinate state. This produces a bilateral valid decreasing chain
\((c_j,c_{j+1},u_j)\) with the four prescribed limits. Choose
\(\lambda_0>0\) and impose \(\lambda_j=u_j\lambda_{j-1}\) in both
directions. All amplitudes are positive. On the far left, ratios are bounded
below by some \(q_->1\), so amplitudes decay geometrically when moving
left. On the far right they are bounded above by some \(q_+<1\), so
amplitudes decay geometrically when moving right. Thus
\(0<\sum\lambda_j^2<\infty\), and a single scaling normalizes the vector.
The first equation in \eqref{eq:stationarypair}, multiplied by
\(\lambda_j/2\), gives exactly \(H(c)\lambda=s\lambda\).
The second equation remains true at every chain step. All conditions of
Definition~\ref{def:carrier} follow.

Conversely, let \((c,\lambda)\in\Cs\). Strict decrease and the positive
left and negative right limits imply that \(\{j:c_j\ge0\}\) is a nonempty
initial segment of \(\ZZ\) bounded above. Its maximum \(k\) is the unique
crossing index. If \(c_k=0\), the representative starts at that zero and
ends at the next negative label, as required by \eqref{eq:matching}.
The full states \(X_j=(c_j,c_{j+1},\lambda_j/\lambda_{j-1})\) tend to
\(p_+\) and \(p_-\). Hence for some finite \(m,n\ge0\), the entire past
of \(X_{k-m}\) lies in \(N_+\) and the entire future of \(X_{k+n}\) lies
in \(N_-\). Local capture in Lemma~\ref{lem:tails} places these states on
the decreasing unstable and stable pieces. Their actual carrier steps give
\(F_s^m(X_{k-m})=X_k\) and \(F_s^n(X_k)=X_{k+n}\), with all validity
conditions retained. Thus \(X_k\in\Ms\).
\end{proof}

\begin{proposition}[A match at the true value]\label{prop:existmatch}
The matching set \(\mathcal M_{\St}\) is nonempty.
\end{proposition}
\begin{proof}
Proposition~\ref{prop:R1} puts \(\St\) in the domain \(I\).
Theorem~\ref{thm:R0} supplies a member of \(\mathcal C_{\St}\).
Apply Proposition~\ref{prop:bridge} at that parameter.
\end{proof}
